\section*{Exercises}
Exercises~\ref{ex:8.1}--\ref{ex:8.4} are core exercises, and Exercises~\ref{ex:8.5} and~\ref{ex:8.6} are transfer exercises. Complete solutions are in Appendix~\ref{app:solb}, under ``Week~8.''
\begin{exercise}\label{ex:8.1}In the Ada--Ben graph of Section~\ref{ch08:sec:model}, Ada accepts tasks $1$ and $2$, and Ben accepts only task~$1$. List all the matchings of this graph, decide which of them are maximal and which are maximum, and conclude that $\{\text{Ada--}1\}$ is the only matching that is maximal but not maximum.\end{exercise}
\begin{exercise}\label{ex:8.2}Students $a$, $b$, $c$ accept the tasks $N(a)=\{1,2\}$, $N(b)=\{1,3\}$, $N(c)=\{2\}$. Start from the matching $\{a\text{--}2,\ b\text{--}1\}$. Use the search of Section~\ref{ch08:sec:search} to find an augmenting path, and write down the matching obtained by flipping along it.\end{exercise}
\begin{exercise}\label{ex:8.3}Three students each accept exactly the tasks $1$ and $2$, as in the example of Section~\ref{ch08:sec:shortage}. Give a shortage certificate, and explain why the fact that every student has two acceptable tasks is not enough for a full assignment. Then show that no smaller group of these students is a shortage set.\end{exercise}
\begin{exercise}\label{ex:8.4}The proof of Lemma~\ref{thm:augment} uses the fact that an augmenting path has no repeated vertex. (a)~Explain where the proof uses this fact. Give an example of a matching and an alternating walk from an unmatched student to an unused task that visits some vertex twice, such that flipping along the walk does not give a matching. (b)~Suppose that some alternating walk leads from an unmatched student to an unused task. Prove that a shortest such walk has no repeated vertex, so that it is an augmenting path.\end{exercise}
\begin{exercise}\label{ex:8.5}A procedure works with matchings in a bipartite graph whose set of students is $L$, and each of its steps enlarges the current matching by exactly one edge, for example by flipping along an augmenting path. Explain why the fact that the matching grows at every step does not, on its own, prove that the procedure stops, and why the finiteness of $L$ supplies the missing ingredient. State an upper bound on the number of steps when the procedure starts from a matching with $k$ edges.\end{exercise}
\begin{exercise}\label{ex:8.6}Suppose that the complete search of Section~\ref{ch08:sec:search}, started at an unmatched student $x$, fails, and let $A$ and $B$ be the sets of students and tasks it reached. (a)~Without looking back at the proof of Proposition~\ref{ch08:prop:shortage}, prove carefully that $N(A)=B$, treating matched and unmatched edges separately, and point out exactly where the completeness of the search is used. (b)~Now suppose instead that the search is stopped early: it has reached no unused task, but some reached student has not been completely inspected. Let $A$ and $B$ again be the sets of students and tasks reached so far. Which of the two containments $B\subseteq N(A)$ and $N(A)\subseteq B$ must still hold? Give an example, different from the one in Section~\ref{ch08:sec:shortage}, in which the other one fails. Then show that whenever it fails, $A$ is not a shortage set.\end{exercise}
